%% brochure-body.tex — JA4proxy Product Overview & Buying Guide
%% Body content for brochure.tex (tufte-handout, A4)
%% Revised for v2.0.0 Go-Native Release

% -----------------------------------------------------------------------------
\section{Stop the Bots That Look Human}
\index{bots!modern evasion}\index{SSL inspection}\index{compliance}
% -----------------------------------------------------------------------------

\newthought{A credential-stuffing campaign looks identical to human login traffic}
at the HTTP layer. The requests are syntactically valid. The user-agent strings are
plausible. Your CDN's bot-score misses it because the attacker rotates IP addresses
faster than reputation feeds update.

\sidenote{The Verizon DBIR consistently shows that credential abuse accounts for
more than 40\% of breaches. Most credential-stuffing toolkits ship with TLS stacks
that produce distinctive, stable fingerprints --- yet traditional defences never
inspect the handshake.}

The usual answer --- TLS inspection and decryption --- is a poor trade. It adds
latency, demands key custody, and creates a compliance exposure that PCI-DSS,
HIPAA, and GDPR all penalise. \ja{} takes the opposite approach: it makes its
decision \emph{before} decryption, reading only the plaintext the client
volunteers in its ClientHello. It never holds a key, never decrypts a byte, and
forwards allowed traffic through unchanged.

% -----------------------------------------------------------------------------
\section{The JA4 Insight: Identity Before Encryption}
\index{JA4!fingerprint}\index{TLS!ClientHello}
% -----------------------------------------------------------------------------

\newthought{Before a single byte of encrypted data is exchanged,} every TLS client
announces its capabilities in cleartext --- TLS version, cipher suites, and
extensions (SNI, ALPN). Because that announcement reflects the client's TLS stack,
the fingerprint is remarkably stable: Sliver 1.5.x produces the same ClientHello
whether it runs from Frankfurt or Tokyo, and swapping the \emph{payload} does
nothing to change it. You cannot patch your way out of your own handshake.

\begin{table}[ht]
  \small
  \begin{tabularx}{\linewidth}{L{2.8cm} L{4.6cm} L{2.0cm}}
    \toprule
    \textbf{Client} & \textbf{JA4 Fingerprint (example)} & \textbf{Disposition} \\
    \midrule
    Chrome 125 (Win) & \fingerprint{t13d1516h2\_8daaf6152771\_02713d6af862} & \badgeallow \\[4pt]
    Firefox 126      & \fingerprint{t13d1915h2\_cd87aab0d655\_5d1a4de8bd19} & \badgeallow \\[4pt]
    Sliver C2        & \fingerprint{t13d2012\_c1e4e3a9b26a\_000000000000} & \badgeblock \\[4pt]
    Cobalt Strike    & \fingerprint{t13d1112\_b4f47bba57b5\_9b5e0e9c9d74} & \badgeblock \\[4pt]
    \bottomrule
  \end{tabularx}
  \caption{Representative JA4 fingerprints. Browser fingerprints carry an \code{h2}
           field; C2 frameworks typically do not.}
\end{table}

% -----------------------------------------------------------------------------
\section{Throughput, Not Trade-offs}
\index{performance}\index{throughput}\index{Go-native}
% -----------------------------------------------------------------------------

\newthought{\ja{} v2.0.0 is a pure Go-native engine} that sustains roughly
\textbf{2,500 connections per second per node}\sidenote{Measured end-to-end on
commodity hardware (Intel i9-9900K) with host networking and the full scoring
pipeline active --- zero errors, zero false positives. Bridged container
networking is lower; deploy with host networking for line-rate throughput.}
--- enough to front a busy site from a single modest box, scaling linearly as
you add nodes.

That headroom comes from a deliberately small core: a fingerprint decision on
the allow path costs about \textbf{272 ns}, so the proxy adds latency your users
will never feel. Layer-by-layer isolation testing proved the security pipeline
itself is never the bottleneck --- your network is.

\begin{table}[ht]
  \small
  \begin{tabularx}{\linewidth}{L{3.5cm} R{2.5cm} R{2.5cm}}
    \toprule
    \textbf{Decision path} & \textbf{Core latency} & \textbf{Heap allocs} \\
    \midrule
    \signalrow{Allowed (Fast Path)}    & \textbf{272 ns} & 1 object \\[3pt]
    \signalrow{Scored (Standard Path)}  & \textbf{2.3 $\mu$s} & 9 objects \\[3pt]
    \signalrow{Adversarial (JA4X)}      & \textbf{6.2 $\mu$s} & 73 objects \\
    \bottomrule
  \end{tabularx}
  \caption{Per-connection pipeline cost, isolated from network I/O
           (Intel i9-9900K). A single static binary, under 12 MB resident.}
\end{table}

% -----------------------------------------------------------------------------
\section{It Cannot Block Your Real Customers}
\index{false positive}\index{fail open}\index{monitor mode}
% -----------------------------------------------------------------------------

\newthought{The expensive failure for a security proxy is a false positive} --- a
real browser turned away. \ja{} is built so that outcome is structurally hard to
reach.

\begin{itemize}
  \item \textbf{Monitor-mode by default.} Out of the box the dial sits at
        \dialzero{}: every connection is scored and logged, nothing is blocked.
        You raise the dial only once you trust the scores.
  \item \textbf{Browser traffic bypasses everything.} \code{h2}/\code{h1} ALPN
        handshakes are allowed before the scorer ever runs --- a real browser
        cannot be blocked, by design.
  \item \textbf{Fail open, always.} If Redis or any enrichment service is
        unreachable, \ja{} allows and logs. A degraded dependency never becomes
        an outage for your users.
\end{itemize}

\sidenote{No decryption means no key custody, no certificate logistics, and no
new audit scope. \ja{} is transparent to both client and backend, and forwards
allowed connections byte-for-byte.}

% -----------------------------------------------------------------------------
\section{Deploys in Minutes}
\index{operations}\index{deployment}\index{v2.0.0}
% -----------------------------------------------------------------------------

\newthought{The engine is one statically-linked Go binary} in a minimal
container image --- no runtime, no interpreter, nothing to patch underneath it.
Bring up the bundled Docker Compose stack, or run the container under the
orchestrator of your choice, and \code{ja4pd} starts in milliseconds, ready for
auto-scaling DMZ deployments. The whole project is open source: clone the v2.0.0
release and \code{make build} the engine and its companion CLI (\code{ja4p}) from
source whenever you prefer.

For high-traffic deployments the engine adds integrated cluster
synchronisation, Zero-Trust Redis hardening over TLS, and PROXY protocol v2 with
TLV-based trust verification. Real-time Prometheus and Grafana metrics give you
granular visibility into the full security waterfall --- without costing you
throughput.

\vspace{\baselineskip}
{\small\sffamily\color{ja4gray}
Source code and technical evaluators' guide:\\
\url{https://github.com/seanpor/JA4proxy}\\[2pt]
JA4 Specification: \url{https://github.com/FoxIO-LLC/ja4}
}
